% ===== wrap-up =====
\begin{frame}{이번 주차 정리}
  \begin{enumerate}\small
    \item \kb{9.1 — 은닉층이 XOR의 벽을 넘는다}: 부등식 모순으로
      단층 퍼셉트론의 한계를 보이고, 은닉층이 ``직선으로 갈리는
      새 공간''을 만든다. 역전파 = 연쇄법칙의 \textbf{재사용}
      (동적 프로그래밍) — 수치 미분으로 구현을 검증한다.
    \item \kb{9.2 — 그래디언트는 ``할인율의 곱''}: 층마다 $\sigma'
      \times W$를 곱해 깊어지면 지수적으로 0(소실)이거나 불어나
      (폭발). 세 라인 방어: \textbf{활성함수(ReLU 계열) + 초기화
      (Xavier/He) + 정규화(BatchNorm)}.
    \item \kb{9.3 — 학습을 실제로 끝내는 도구}: dropout(공적응 억제) ·
      BatchNorm(분포 안정, 큰 lr 가능) · 가중치 감쇠(AdamW) ·
      학습률 스케줄링(초반 크게, 후반 작게, warmup).
  \end{enumerate}
  \vskip0.8em
  \begin{block}{다음 주 --- Week 10. 합성곱 신경망 (CNN)}
    이번 장의 ``일반부''(역전파, 활성함수, 초기화, 정규화) 위에,
    \textbf{이미지의 특성에 맞는 ``특수부''} — 합성곱층 — 를 올린다.
    이번 장에서 세운 역전파·활성함수·초기화·정규화·스케줄링은
    그 뒤의 모든 장(CNN$\to$RNN$\to$트랜스포머$\to$LLM)에서 다시
    쓰이는 도구다.
  \end{block}
\end{frame}
